\documentclass[11pt]{article}
\usepackage{ctex}
\usepackage{amsmath,amssymb}
\usepackage{geometry}
\geometry{a4paper,margin=25mm}
\title{直接求解 Euler--Lagrange 微分方程的详细步骤}
\author{}
\date{}

\begin{document}
\maketitle

\section*{问题}
从变分问题得到的欧拉--拉格朗日方程（引入拉格朗日乘子 $\lambda$ 后）为
\[
\frac{d}{dx}\!\left(\frac{y'}{\sqrt{1+y'^2}}\right)=\lambda,
\qquad x\in[x_0,x_1].
\tag{1}
\]
要求对该微分方程作解析求解，并把常数关系与几何意义写明。

\section*{第一步：做替换，化为一阶方程}
令辅助函数
\[
p(x):=\frac{y'(x)}{\sqrt{1+y'(x)^2}}.
\]
由定义可见 $p(x)\in(-1,1)$ 且 $p$ 与 $y'$ 单射可互化。方程 (1) 变为
\[
p'(x)=\lambda.
\]
对 $x$ 积分得到（$C$ 为积分常数）
\[
\boxed{\,p(x)=\lambda x + C\,.}
\tag{2}
\]

\section*{第二步：由 $p$ 恢复 $y'$}
由 $p$ 的定义，
\[
p=\frac{y'}{\sqrt{1+y'^2}}
\]
两边平方并解出 $y'^2$：
\[
p^2=\frac{y'^2}{1+y'^2}
\quad\Longrightarrow\quad
y'^2=\frac{p^2}{1-p^2}.
\]
因此取正/负号（由问题的单调性或端点条件决定）
\[
\boxed{\,y'(x)=\frac{p(x)}{\sqrt{1-p(x)^2}}\,.}
\tag{3}
\]
将 (2) 代入 (3)：
\[
\boxed{\,y'(x)=\frac{\lambda x + C}{\sqrt{1-(\lambda x + C)^2}}\,.}
\tag{4}
\]

\section*{第三步：积分得到 $y(x)$}
做换元 $u=\lambda x + C$，则 $du=\lambda\,dx$，从而
\[
dy = y'\,dx = \frac{u}{\sqrt{1-u^2}}\cdot\frac{du}{\lambda}.
\]
对两边积分：
\[
y(x)=\int \frac{u}{\lambda\sqrt{1-u^2}}\,du + D
= -\frac{1}{\lambda}\sqrt{1-u^2} + D,
\]
其中 $D$ 为积分常数。把 $u=\lambda x + C$ 代回，得到显式形式
\[
\boxed{\,y(x) = -\frac{1}{\lambda}\sqrt{1-(\lambda x + C)^2} + D\,.}
\tag{5}
\]

\section*{第四步：化为圆的标准方程}
对 (5) 两边适当代数变形可得到圆的方程。先将等式改写为
\[
\lambda\bigl(D-y(x)\bigr)=\sqrt{1-(\lambda x + C)^2}.
\]
两边平方：
\[
\lambda^2 (D-y)^2 = 1 - (\lambda x + C)^2.
\]
移项并整理：
\[
(\lambda x + C)^2 + \lambda^2 (D-y)^2 = 1.
\]
令
\[
R=\frac{1}{\lambda},\qquad x_c=-\frac{C}{\lambda},\qquad y_c=D,
\]
则上式等价于
\[
\boxed{(x-x_c)^2 + (y-y_c)^2 = R^2,}
\]
即解为以 $(x_c,y_c)$ 为圆心、半径 $R$ 的圆的上（或下）半弧。这里 $\lambda=0$ 时需单独讨论（见下）。

\section*{第五步：曲率解释与特例}
\begin{itemize}
  \item 从方程 (1) 可直接得到
  \[
  \frac{y''}{(1+y'^2)^{3/2}}=\lambda,
  \]
  而左边正是平面曲线 $y(x)$ 的曲率 $\kappa(x)$。因此 $\kappa(x)=\lambda$ 为常数，故曲线为圆（当 $\lambda\neq0$），当 $\lambda=0$ 时曲率为 $0$，曲线退化为直线（$y'=\text{const}$）。
  \item $\lambda$ 的正负决定圆弧的凸性：$\lambda>0$ 对应向上凸（圆心在曲线的上方），$\lambda<0$ 对应向下凸。
\end{itemize}

\section*{第六步：由边界条件与面积约束确定常数}
若给定两个端点条件
\[
y(x_0)=y_0,\qquad y(x_1)=y_1,
\]
将它们代入圆方程
\[
(x_0-x_c)^2 + (y_0-y_c)^2 = R^2,
\qquad
(x_1-x_c)^2 + (y_1-y_c)^2 = R^2,
\]
可得到关于未知量 $(x_c,y_c,R)$ 的两条方程；再若有面积约束（例如
\[
\int_{x_0}^{x_1} y(x)\,dx = A_0
\]
），则该积分表达式（将 $y(x)$ 以 (5) 形式代入并积分）将给出第三个关于 $(x_c,y_c,R)$ 的方程，从而可确定三个未知常数。对于对称情形（例如本题中 $x_0=0,x_1=1,y_0=y_1=0$），可利用对称性令 $x_c=\tfrac{1}{2}$，简化求解：
\[
R^2=\Bigl(\tfrac{1}{2}\Bigr)^2 + h^2,\quad y_c=h,
\]
并由圆段面积公式
\[
A_0 = R^2\Bigl(\theta - \sin\theta\cos\theta\Bigr),\qquad \sin\theta=\frac{1}{2R},
\]
求出 $R$（从而 $\lambda=1/R$）和 $h$。
\medskip

\noindent 以上给出了解常数确定的步骤；在一般非对称情况下，代入边界条件与面积约束会得到关于 $\lambda,C,D$ 的代数或超越方程，通常需数值方法求解。

\section*{第七步：处理符号与退化情况}
\begin{itemize}
  \item 在 (3) 中，选择 $y'$ 的正负号应由问题中曲线的单调性或端点斜率决定；若 $y$ 在区间上单调增加则取 $+$，单调减少则取 $-$。
  \item 退化情况 $\lambda=0$：从 (1) 有 $\dfrac{d}{dx}\!\bigl(\dfrac{y'}{\sqrt{1+y'^2}}\bigr)=0$，即 $\dfrac{y'}{\sqrt{1+y'^2}}=C$，故 $y'=$ 常数，解为直线 $y=ax+b$。
\end{itemize}

\section*{结论}
对 Euler--Lagrange 给出的方程作直接求解，经过一次积分与代换，可以得到显式的 $y'(x)$ 及 $y(x)$ 形式，并可证明通解为圆（或直线）的方程：
\[
(x-x_c)^2 + (y-y_c)^2 = R^2,
\qquad R=\frac{1}{\lambda},\ x_c=-\frac{C}{\lambda},\ y_c=D.
\]

\vfill
\noindent ----- 结束 -----

\end{document}
